\section{Lesson 2: Why Am I Supposed to Believe It?}

Restating a sentence accurately and fixing its terms clears the grammar, yet leaves the decisive question untouched: \emph{why am I supposed to accept this?} Within ten lines of a single article, Aquinas may state a definition, appeal to observation, cite an argument proved two questions earlier, lay down a causal principle, quote Aristotle, and quote St.\ Paul---all in the same brisk declarative voice. Unless you can sort out what supports each premise, every sentence blurs into one uniform wall of assertion.

Once claim and meaning are fixed, the third question enters: \textbf{Warrant}. Two rules govern it. First, warrant attaches to \emph{premises one at a time}, never to whole arguments. Second, classical authors rarely announce which kind of support each sentence carries, so the reader must judge explicitly and acknowledge genuine ambiguity where a sentence admits two readings. Hold fast also to our standing rule against anachronism from Lesson 0: the eight kinds of warrant below are our modern sorting list, never Aristotle's six conditions on demonstrative premises. Two of our eight---\emph{revelation} and \emph{provenance only}---have no place in Aristotle's theory of premises at all.

\subsection{The Eight Kinds of Warrant, Seen on Real Sentences}

Watch all eight kinds of warrant at work on real sentences from Aquinas's \emph{Summa Theologiae} (1920 English Dominican translation):

\begin{enumerate}
\item \textbf{Definition.} ``For a body is that which has the three dimensions'' (\emph{ST} I, q.\,3, a.\,1, Obj.\ 1). True by what the term means; asserts nothing yet about what exists.
\item \textbf{Logical consequence.} Given that no body moves unless put in motion and that God is the unmoved First Mover, ``Therefore it is clear that God is not a body'' (\emph{ST} I, q.\,3, a.\,1, \emph{corpus}). Support is purely the validity of the inference from premises already laid down.
\item \textbf{Experience.} ``It is certain, and evident to our senses, that in the world some things are in motion'' (\emph{ST} I, q.\,2, a.\,3, \emph{corpus}). Support rests on plain, everyday encounter.
\item \textbf{Prior result.} ``Now it has been shown above (Q.\,3, A.\,4) when treating of the divine simplicity that God is the essentially self-subsisting Being'' (\emph{ST} I, q.\,44, a.\,1, \emph{corpus}). A promissory warrant worth exactly what the cited argument proved.
\item \textbf{Metaphysical principle.} ``There is no case known (neither is it, indeed, possible) in which a thing is found to be the efficient cause of itself; for so it would be prior to itself, which is impossible'' (\emph{ST} I, q.\,2, a.\,3, \emph{corpus}); or ``whatever is found in anything by participation, must be caused in it by that to which it belongs essentially'' (\emph{ST} I, q.\,44, a.\,1, \emph{corpus}). Neither a bare definition nor a direct observation, but a substantive claim about reality offered as a starting point.
\item \textbf{Hypothesis.} ``Further, if the world always was, the consequence is that infinite days preceded this present day'' (\emph{ST} I, q.\,46, a.\,2, Obj.\ 6). Assumed provisionally to test what follows; in wider philosophy, \emph{hypothesis} also names inference to the best available explanation---legitimate reasoning, yet not Aristotelian demonstration.
\item \textbf{Revelation.} ``By faith alone do we hold, and by no demonstration can it be proved, that the world did not always exist'' (\emph{ST} I, q.\,46, a.\,2, \emph{corpus}). Held on the authority of God revealing in Scripture and the Creed.
\item \textbf{Provenance only.} A bare citation standing without argument on the page, such as ``as Damascene says (\emph{De Fide Orth.}\ i, 1, 3), `the knowledge of God is naturally implanted in all''' (\emph{ST} I, q.\,2, a.\,1, Obj.\ 1).
\end{enumerate}

Hold one sentence fixed about that eighth item: \textbf{provenance only is a diagnosis, not a warrant.} Recording ``provenance only'' never refutes a proposition; it may be true and proved three pages later. It records honestly that on \emph{this} page you have been handed a name rather than a reason.

These eight kinds are not ranked in respectability: a definition is not a weak warrant, merely a warrant for a conceptual claim. Single sentences also carry \emph{mixed warrant} (partly definitional, partly substantive) or \emph{layered warrant} (a prior result resting on earlier premises still). Where a text leaves a premise genuinely underdetermined between two kinds, note both readings instead of guessing.

\subsection{Provenance, Status, and Warrant---and the Name-Removal Test}

Mistaking who said a sentence for why it is true is the commonest error in reading classical philosophy. Keep three questions strictly apart:

\begin{itemize}
\item \textbf{Provenance:} Who asserts this proposition, and where? A historical fact settled by citation.
\item \textbf{Status:} What role does the sentence play---definition, objection, dialectical argument, demonstration claimed, demonstration achieved, theological conclusion, or authorial judgment? \emph{Demonstration claimed} is not \emph{demonstration achieved}: calling an argument a demonstration is a fact about the text; whether it succeeds as one is a further question.
\item \textbf{Warrant:} Why should the reader believe the proposition is true? A philosophical question settled by reasons, not prestige.
\end{itemize}

Hold these four distinct sentences apart: ``Aquinas holds P'' (provenance); ``Aquinas argues for P here'' (status of a passage); ``Aquinas cites Aristotle for P'' (layered provenance); and ``Aquinas claims to demonstrate P'' (claimed status, not yet achieved proof).

To detect when citation is impersonating argument, run the \textbf{name-removal test}: \emph{strike every proper name and citation from the passage, and ask whether any premise lost philosophical support.} If the reasoning stands intact, the names were doing legitimate scholarly work---crediting authors and situating a debate. If a premise collapses when ``Aristotle says'' or ``Aquinas holds'' vanishes, the name was carrying weight it cannot bear, exposing \emph{provenance only}. This rule is not hostility to great teachers. Aristotle, Augustine, and Aquinas merit sustained attention because of the depth of their thought, yet philosophical claims stand or fall by their warrant---which is precisely why medieval doctors disputed authorities openly rather than ending inquiries by quotation.

\subsection{Scrutinizing Prior Results and Appeals to Experience}

Two genuine warrants---\textbf{prior result} and \textbf{experience}---require close inspection because both can smuggle unexamined steps past a hurried reader.

A \textbf{prior result} is a promissory note that must be checked at the cited address. Look at a concrete case where the referenced article establishes less than the citing sentence claims. In \emph{Summa Theologiae} I, q.\,3, a.\,2, \emph{corpus}, arguing that God contains no matter, Aquinas writes: ``First, because matter is in potentiality. But we have shown (Q.\,2, A.\,3) that God is pure act, without any potentiality.'' Turn back to \emph{ST} I, q.\,2, a.\,3, \emph{corpus}. The First Way concludes to ``a first mover, put in motion by no other; and this everyone understands to be God.'' It does \emph{not} yet argue that this first mover lacks every kind of potentiality. That broader exclusion of all potentiality is actually argued in \emph{ST} I, q.\,3, a.\,1, \emph{corpus} (moving from ``First Being'' to ``in act, and in no way in potentiality''). The compressed cross-reference to q.\,2, a.\,3 leaps over the intervening step in q.\,3, a.\,1. Checking the promissory note does not refute Aquinas; it shows you which bridge actually carries the load.

\textbf{Experience} demands the same scrutiny. Ask whether a sentence reports plain observation or already clothes observation in the vocabulary the argument needs. Compare two lines from the First Way (\emph{ST} I, q.\,2, a.\,3, \emph{corpus}): ``It is certain, and evident to our senses, that in the world some things are in motion'' against ``Whatever is moved is moved by another'' (or ``we plainly see things passing from potentiality to actuality''). That things change---iron heats, water freezes, leaves fall---is plain observation. That every change is a passive actualization by an external mover is \emph{not} seen by the eye; it is a metaphysical principle dressed as an observation. Even the Latin passive verb \emph{moveri} (``to be moved'') gently nudges a reader toward ``moved by another'' before any proof arrives. Separating observation from theoretical description prepares you to locate the real philosophical burden in Lesson 3.

\subsection{Revealed Premises, Theology, and Natural Reason}

Reading Catholic masters such as Augustine or Aquinas requires separating four kinds of proposition that routinely stand side by side:

\begin{enumerate}
\item \textbf{A revealed premise:} held on divine revelation and presented by the author as such.
\item \textbf{Theological reasoning from revealed premises:} valid deduction from revealed starting points, carrying the revealed standing of its premises.
\item \textbf{A philosophical claim accessible to natural reason:} offered as establishable without appeal to revelation, and therefore answerable to natural argument.
\item \textbf{Authorial judgment:} the author's considered opinion on a predecessor, an interpretation, or the weight of an argument.
\end{enumerate}

Keeping these four distinct prevents \textbf{two symmetrical errors}, both failures of bookkeeping rather than of piety. \textbf{Theology in philosophical clothing} lets a premise known only by faith pass silently as a result of natural reason; the remedy is restoring its label and asking what survives without it. \textbf{Philosophy defended by authority} presents a claim as naturally demonstrable and then retreats behind saintly or ecclesial prestige when pressed; a claim advanced on natural reason accepts natural reason's standards.

This discipline is Aquinas's own. In \emph{Summa Theologiae} I, q.\,1, a.\,2, \emph{corpus}, he separates philosophical sciences, proceeding from principles known by the natural light of intelligence, from \emph{sacra doctrina} (sacred doctrine), proceeding from principles established by the light of a higher science---the science of God and the blessed---revealed by God. In \emph{Summa contra Gentiles} I, c.\,3, he names his \emph{duplex veritatis modus} (``twofold mode of truth'') about God: truths surpassing human reason (such as the Trinity) against truths reachable by natural reason (such as that God exists and is one). In \emph{ST} I, q.\,2, a.\,2, \emph{ad 1}, he calls those naturally accessible truths \emph{praeambula fidei} (``preambles to the articles''), noting that faith presupposes natural knowledge as grace presupposes nature, while adding that anyone unable to follow a demonstration may rightly receive those same truths by faith. Keeping faith and natural reason labelled lets a believer see clearly what philosophy actually proves.

\subsection{Worked Passage 1: A Proposition Whose Warrant Is Revelation}

Read \emph{Summa Theologiae} I, q.\,46, a.\,2, \emph{corpus} (``Whether It Is an Article of Faith That the World Began?'', 1920 English Dominican translation; Leonine Latin: \emph{mundum non semper fuisse, sola fide tenetur, et demonstrative probari non potest}), alongside \emph{SCG} I, c.\,3 and \emph{ST} I, q.\,2, a.\,2, \emph{ad 1}:

\begin{quote}
By faith alone do we hold, and by no demonstration can it be proved, that the world did not always exist, as was said above of the mystery of the Trinity (Q.\,32, A.\,1). The reason of this is that the newness of the world cannot be demonstrated on the part of the world itself. For the principle of demonstration is the essence of a thing. Now everything according to its species is abstracted from ``here'' and ``now''; whence it is said that universals are everywhere and always. Hence it cannot be demonstrated that man, or heaven, or a stone were not always. Likewise neither can it be demonstrated on the part of the efficient cause, which acts by will. For the will of God cannot be investigated by reason, except as regards those things which God must will of necessity; and what He wills about creatures is not among these, as was said above (Q.\,19, A.\,3). But the divine will can be manifested by revelation, on which faith rests. Hence that the world began to exist is an object of faith, but not of demonstration or science. And it is useful to consider this, lest anyone, presuming to demonstrate what is of faith, should bring forward reasons that are not cogent, so as to give occasion to unbelievers to laugh, thinking that on such grounds we believe things that are of faith.
\end{quote}

We carry this passage through \textbf{Claim}, \textbf{Meaning}, and \textbf{Warrant} (Burden waits for Lesson 3), with no symbolic notation: what is at stake is the standing of propositions, not formal scope.

\subsubsection*{Claim}
As Aquinas writes (\emph{ST} I, q.\,46, a.\,2, \emph{corpus}): we hold solely by faith, and cannot prove by demonstration, that the world did not always exist; neither created essences nor God's free will allow natural reason to demonstrate a temporal beginning, though divine revelation manifests it.

\emph{Plain-language reading.} That the universe began in time rather than existing forever is known strictly by divine revelation received in faith. Natural reason cannot prove a temporal beginning either from the natures of creatures (which say nothing about particular dates) or from God's will (since God creates freely rather than of necessity).

\subsubsection*{Meaning}
\emph{Object of faith} (or article of faith) means a truth held directly on divine revelation because it exceeds natural demonstration---unlike the \emph{praeambula fidei} of \emph{ST} I, q.\,2, a.\,2, \emph{ad 1}, placing the world's temporal beginning in the first class of \emph{SCG} I, c.\,3's \emph{duplex veritatis modus}. \emph{Demonstration} means strict Aristotelian proof from essence or necessary cause. \emph{Species} and \emph{universal} name what a thing is in kind, abstracted from ``here'' and ``now.'' \emph{Will of necessity} names what God must will (His own goodness), contrasted with creatures, which He wills freely.

\subsubsection*{Warrant}
The paragraph cleanly separates two claims with different warrants. First, the proposition \emph{that the world did not always exist} rests explicitly on \textbf{revelation} (supported in the \emph{sed contra} by Genesis 1:1 and the Creed). Because Aquinas never claims a philosophical proof for it, demanding one misreads the text; conversely, disguising it as natural reason falsifies Aquinas's explicit teaching.

Second, Aquinas's philosophical claim \emph{that a temporal beginning cannot be demonstrated} rests on natural reason, braiding three warrants: \textbf{prior results} pointing to q.\,32, a.\,1 (the Trinity) and q.\,19, a.\,3 (God wills creatures freely); a \textbf{metaphysical principle} paired with the Aristotelian \textbf{definition} of demonstration (demonstration proceeds from essence, and species abstract from ``here'' and ``now''); and \textbf{logical consequence} concluding that no essence-based or cause-based proof can establish a first moment.

Declining a premise warranted by \textbf{revelation} incurs neither a logical contradiction nor an explanatory cost in the philosophical sense; someone who declines revelation simply does not share that starting point. Keeping revealed premises plainly labelled honours both faith and natural reason, vindicating Aquinas's own closing warning against weak ``proofs'' that provoke unbelievers to laugh and make them suppose Christians hold articles of faith on such grounds.

\subsection{Worked Passage 2: Provenance Mistaken for Warrant}

Our second worked passage shows how mistaking citation for warrant turns an article upside down. Read four excerpts from \emph{Summa Theologiae} I, q.\,2, a.\,1 (``Whether the Existence of God Is Self-Evident?'', 1920 English Dominican translation), watching \emph{where} each stands in the \emph{disputatio}. First, Objection 1:

\begin{quote}
Objection 1: It seems that the existence of God is self-evident. Now those things are said to be self-evident to us the knowledge of which is naturally implanted in us, as we can see in regard to first principles. But as Damascene says (De Fide Orth.\ i, 1, 3), ``the knowledge of God is naturally implanted in all.'' Therefore the existence of God is self-evident.
\end{quote}

Next, the \emph{sed contra}, citing Aristotle and the Psalms against self-evidence:

\begin{quote}
\emph{On the contrary,} No one can mentally admit the opposite of what is self-evident; as the Philosopher (Metaph.\ iv, lect.\ vi) states concerning the first principles of demonstration. But the opposite of the proposition ``God is'' can be mentally admitted: ``The fool said in his heart, There is no God'' (Ps.\ 52:1). Therefore, that God exists is not self-evident.
\end{quote}

Then the conclusion of Aquinas's \emph{corpus}, denying self-evidence \emph{to us}:

\begin{quote}
Therefore I say that this proposition, ``God exists,'' of itself is self-evident, for the predicate is the same as the subject, because God is His own existence as will be hereafter shown (Q.\,3, Art.\ 4). Now because we do not know the essence of God, the proposition is not self-evident to us; but needs to be demonstrated by things that are more known to us, though less known in their nature---namely, by effects.
\end{quote}

Finally, Aquinas's Reply to Objection 1, qualifying St.\ John of Damascus:

\begin{quote}
Reply Obj.\ 1: To know that God exists in a general and confused way is implanted in us by nature, inasmuch as God is man's beatitude. For man naturally desires happiness, and what is naturally desired by man must be naturally known to him. This, however, is not to know absolutely that God exists; just as to know that someone is approaching is not the same as to know that Peter is approaching, even though it is Peter who is approaching; for many there are who imagine that man's perfect good which is happiness, consists in riches, and others in pleasures, and others in something else.
\end{quote}

Read this beside Aquinas's own hierarchy of authorities in \emph{Summa Theologiae} I, q.\,1, a.\,8, \emph{ad 2} (same translation), where he writes that sacred doctrine uses philosophers as ``extrinsic and probable arguments,'' the doctors of the Church properly yet ``merely as probable,'' and canonical Scripture ``as an incontrovertible proof,'' quoting St.\ Augustine's letter to Jerome (\emph{Epist.\ ad Hieron.}\ xix, 1):

\begin{quote}
Only those books of Scripture which are called canonical have I learned to hold in such honor as to believe their authors have not erred in any way in writing them. But other authors I so read as not to deem everything in their works to be true, merely on account of their having so thought and written, whatever may have been their holiness and learning.
\end{quote}

Again, we stop at \textbf{Warrant} with no notation; what matters is what each citation is doing.

\subsubsection*{Claim}
As Aquinas determines in \emph{ST} I, q.\,2, a.\,1, \emph{corpus} and \emph{ad 1}: ``God exists'' is self-evident in itself, because God is His own existence, yet \emph{not} self-evident to us, because we do not know God's essence; Damascene's statement indicates only a general, confused natural desire for happiness, not immediate knowledge of God.

\emph{Plain-language reading.} Although ``God exists'' is necessarily true by what God is in Himself, we do not directly grasp God's nature and therefore cannot see immediately that God exists; we must demonstrate it from effects. Damascene's saying means only that humans naturally desire happiness, not that they immediately recognize God as its object.

\subsubsection*{Meaning}
\emph{Self-evident} (\emph{per se notum}) means that the predicate belongs inside the essence of the subject (as ``animal'' belongs inside ``man''). \emph{In itself} against \emph{to us} (\emph{quoad se} against \emph{quoad nos}) divides objective inclusion in the subject's essence from human grasp of that essence. \emph{General and confused knowledge} means apprehending something indeterminately without identifying it specifically---knowing someone approaches without knowing it is Peter.

\subsubsection*{Warrant}
Run the name-removal test across the article. Objection 1's appeal to St.\ John of Damascus carries \textbf{provenance only} as philosophical proof: strike Damascene's name, and no reason remains to think direct knowledge of God is implanted in every mind. Worse, the quotation sits inside \emph{Objection 1}---the position Aquinas rejects! A reader who treats citation as warrant reads the entire article backwards. (The \emph{sed contra} cites Aristotle and Psalm 52:1 to set up the opposing claim that people mentally deny ``God is,'' leaving the \emph{corpus} to argue the solution.)

Inside the \emph{corpus} and \emph{ad 1}, strip every name and the argument stands on four clear warrants: a \textbf{definition} of \emph{per se notum}; a \textbf{prior result} (actually a forward promissory note to q.\,3, a.\,4) that God is His own existence; \textbf{experience} that we do not directly inspect God's essence and that people seek happiness in riches or pleasure; and \textbf{logical consequence} concluding that ``God exists'' is not self-evident to us.

The outcome is crisp: the citation establishes only that Damascene wrote those words. Aquinas himself---echoing Augustine in \emph{ST} I, q.\,1, a.\,8, \emph{ad 2}---is the one distinguishing the uses of authority; the confusion belongs solely to a reader who mistakes provenance for warrant.

\subsection{Exercises}

\subsubsection*{Exercise 2.1}
Assign one of the eight kinds of warrant to each numbered sentence from \emph{Summa Theologiae} I, q.\,4, a.\,1, \emph{corpus} (1920 English Dominican translation), explaining your reasoning:

\begin{enumerate}
\item ``As the Philosopher relates (Metaph.\ xii), some ancient philosophers, namely, the Pythagoreans and Leucippus, did not predicate `best' and `most perfect' of the first principle.''
\item ``we call that perfect which lacks nothing of the mode of its perfection.''
\item ``For just as matter, as such, is merely potential, an agent, as such, is in the state of actuality.''
\item ``Hence, the first active principle must needs be most actual, and therefore most perfect.''
\end{enumerate}

\subsubsection*{Exercise 2.2}
Read this excerpt from \emph{Summa Theologiae} I, q.\,44, a.\,1, \emph{corpus} (1920 English Dominican translation):

\begin{quote}
It must be said that every being in any way existing is from God. For whatever is found in anything by participation, must be caused in it by that to which it belongs essentially, as iron becomes ignited by fire. Now it has been shown above (Q.\,3, A.\,4) when treating of the divine simplicity that God is the essentially self-subsisting Being; and also it was shown (Q.\,11, AA.\,3, 4) that subsisting being must be one; as, if whiteness were self-subsisting, it would be one, since whiteness is multiplied by its recipients. Therefore all beings apart from God are not their own being, but are beings by participation. Therefore it must be that all things which are diversified by the diverse participation of being, so as to be more or less perfect, are caused by one First Being, Who possesses being most perfectly.

Hence Plato said (Parmen.\ xxvi) that unity must come before multitude; and Aristotle said (Metaph.\ ii, text 4) that whatever is greatest in being and greatest in truth, is the cause of every being and of every truth; just as whatever is the greatest in heat is the cause of all heat.
\end{quote}

Run the name-removal test: strike Plato and Aristotle, and check the backward references to Q.\,3 and Q.\,11. Separate what rests on a \textbf{prior result}, what stands as a \textbf{metaphysical principle}, and what role Plato and Aristotle play.

\subsubsection*{Exercise 2.3}
Compare two adjacent sentences from the First Way (\emph{ST} I, q.\,2, a.\,3, \emph{corpus}, 1920 English Dominican translation):
\begin{itemize}
\item \emph{Sentence A:} ``For motion is nothing else than the reduction of something from potentiality to actuality.''
\item \emph{Sentence B:} ``But nothing can be reduced from potentiality to actuality, except by something in a state of actuality.''
\end{itemize}
Explain how their warrants differ, and what error occurs if someone treats B as a definition or A as raw sensory experience.

\subsubsection*{Exercise 2.4}
Read the \emph{sed contra} and first argument of the \emph{corpus} of \emph{ST} I, q.\,3, a.\,1 (1920 English Dominican translation):

\begin{quote}
\emph{On the contrary,} It is written in the Gospel of St.\ John (John 4:24): ``God is a spirit.''

\emph{I answer that,} It is absolutely true that God is not a body; and this can be shown in three ways. First, because no body is in motion unless it be put in motion, as is evident from induction. Now it has been already proved (Q.\,2, A.\,3), that God is the First Mover, and is Himself unmoved. Therefore it is clear that God is not a body.
\end{quote}

Separate the warrant of the \emph{sed contra} from the warrants in the \emph{corpus}, say whether the \emph{corpus} borrows from John 4:24, and scrutinize the appeal to ``induction.''

\subsubsection*{Exercise 2.5}
Read this unannotated paragraph---the \emph{corpus} of \emph{Summa Theologiae} I, q.\,1, a.\,2 (1920 English Dominican translation):

\begin{quote}
Sacred doctrine is a science. We must bear in mind that there are two kinds of sciences. There are some which proceed from a principle known by the natural light of intelligence, such as arithmetic and geometry and the like. There are some which proceed from principles known by the light of a higher science: thus the science of perspective proceeds from principles established by geometry, and music from principles established by arithmetic. So it is that sacred doctrine is a science because it proceeds from principles established by the light of a higher science, namely, the science of God and the blessed. Hence, just as the musician accepts on authority the principles taught him by the mathematician, so sacred science is established on principles revealed by God.
\end{quote}

With nothing supplied, work this passage through \textbf{Claim} (with \textbf{Plain-language reading}), \textbf{Meaning}, and \textbf{Warrant}.

\subsection{Solutions}

\subsubsection*{Solution to Exercise 2.1}
Sentence 1 is a historical report; as support for any thesis, it is \textbf{provenance only}. Sentence 2 is a \textbf{definition}, fixing what ``perfect'' means. Sentence 3 functions primarily as a \textbf{metaphysical principle}---with a reasonable case for \emph{mixed warrant}, since Aristotelian ``matter'' and ``agent'' are partly defined through potency and act, yet claiming that real agents act strictly through actuality carries substantive weight. Sentence 4 is a \textbf{logical consequence} of sentences 2 and 3 together with the intermediate claim in q.\,4, a.\,1 that perfection corresponds to actuality.

\subsubsection*{Solution to Exercise 2.2}
Strike Plato and Aristotle: the first paragraph's argument loses nothing. Their names serve as post-argument scholarly corroboration (\textbf{provenance only} if read alone). Inside the first paragraph, ``whatever is found in anything by participation, must be caused in it by that to which it belongs essentially'' is a \textbf{metaphysical principle} (illustrated by fire igniting iron, an appeal to \textbf{experience}); the claims cited to Q.\,3, A.\,4 and Q.\,11, AA.\,3, 4 are \textbf{prior results}---promissory notes whose force depends on those earlier articles; and the two ``Therefore'' sentences follow by \textbf{logical consequence}. The contrast is sharp: citing Q.\,3 and Q.\,11 points to inspectable arguments, not bare prestige.

\subsubsection*{Solution to Exercise 2.3}
Sentence A functions as a \textbf{definition}, fixing what ``motion'' means in the First Way's vocabulary. Sentence B functions as a \textbf{metaphysical principle}, asserting a substantive causal requirement governing every transition from potency to act. Mistaking B for a definition makes the causal claim look trivially verbal; mistaking A for raw sensory experience overlooks that our eyes see logs burning or balls rolling, not ``reductions from potentiality to actuality.'' (Whether A embeds an ontology inside a definition is Lesson 3's question.)

\subsubsection*{Solution to Exercise 2.4}
The \emph{sed contra} rests on \textbf{revelation} (John 4:24). The \emph{corpus} keeps revelation strictly separate and argues by natural reason: ``no body is in motion unless it be put in motion, as is evident from induction'' appeals to \textbf{experience}; ``God is the First Mover, and is Himself unmoved'' cites a \textbf{prior result} (Q.\,2, A.\,3); and ``God is not a body'' follows as a \textbf{logical consequence}. Scrutinizing the appeal to experience catches a live philosophical pressure: what medieval physics took as evident by induction from carts and stones is challenged by modern inertial motion, showing how easily theoretical physics wears the dress of pure observation.

\subsubsection*{Solution to Exercise 2.5}
\textbf{Claim.} As Aquinas writes (\emph{ST} I, q.\,1, a.\,2, \emph{corpus}): sacred doctrine is a science because sciences are of two kinds---some starting from principles known by natural intelligence, others from principles established by a higher science---and sacred doctrine proceeds from principles revealed by God from the higher science of God and the blessed.

\emph{Plain-language reading.} Sacred doctrine qualifies as a science in the second of two senses: just as perspective takes its starting principles from geometry and music takes its principles from arithmetic, sacred doctrine does not begin from naturally self-evident premises but receives its starting truths on authority from God's own knowledge and that of the blessed.

\textbf{Meaning.} \textbf{Science} (\emph{scientia}) means ordered Aristotelian knowledge deduced from firm principles, not modern empirical science. \textbf{Higher science} names a discipline whose conclusions serve as starting points for a subordinate discipline (Aquinas illustrates this concretely with perspective and music without adding a further technical label). \textbf{Sacred doctrine} (\emph{sacra doctrina}) means theology grounded in divine revelation.

\textbf{Warrant.} The twofold division of sciences functions as a \textbf{definition} classifying \emph{scientia}, illustrated from \textbf{experience} of actual disciplines (arithmetic, geometry, perspective, music). That sacred doctrine receives its starting principles from ``the science of God and the blessed'' revealed by God rests on \textbf{revelation}. That sacred doctrine therefore qualifies as a science under the second class follows as a \textbf{logical consequence}. Aquinas neither claims to prove revelation by natural reason nor reduces theology to unreasoned sayings; granting its revealed starting points, he shows why sacred doctrine has the ordered structure of a science.
